\documentclass[11pt]{article}
\usepackage{amsmath,amssymb}
\usepackage[hidelinks]{hyperref}
\title{When a live duty makes a forecast unidentifiable}
\date{}
\begin{document}
\maketitle

Two routing mechanisms require opposite preparation before a fault,
but neither can be identified from the available logs. The issue is
not the number of runs or the accuracy of individual readings.
An existing duty prevents observing a complete sequence in any one
run. The model below gives an exact error--cost frontier when that
experiment becomes feasible. In this model adaptive selection of
local readings cannot improve on randomized advance reservation.

\section{Source process, decisions and available interventions}
Consider a sequence of \(n\ge2\) handoff stages at which two outwardly
indistinguishable carriers can retain or exchange lanes. Stage \(i\)
has actual exchange bit \(X_i\). A fixed but unknown routing mode
\(s\in\{0,1\}\) enforces \(X_1\oplus\cdots\oplus X_n=s\); conditional
on that mode each run is an independent uniform choice of an exchange
sequence with that parity. This is a stipulated source family, not
an estimate from field traces. The final carrier on the initially
loaded carrier's continuation determines which of two positions needs
a movable reserve. There is one reserve. Install it before the
future fault, when relocation would no longer meet the deadline.
A wrong placement invokes a safe slower fallback with additional
service loss \(\ell>0\). Both preparations remain physically safe.
The unknown decision-relevant quantity is the mode, equivalently
the difference of expected losses of the two preparatory actions:
\[
  \mathbb E_s[\ell 1_{\{0\ne s\}}]-\mathbb E_s[\ell 1_{\{1\ne s\}}].
\]
It has opposite signs in the two laws.

A diagnostic station can watch one exchange in its live window;
it reads the lane transition rather than a summary of the end-to-end
carrier identity. A reading at stage \(i\) using option \(u\) costs
\(c_i(u)\ge0\) and returns \(Y_i=X_i\oplus Z_i\) with
\(\mathbb P(Z_i=1)=p_i(u)\in[0,1/2]\). The acquisition noises are
fresh and independent of each other and the source, conditional on
past history and the chosen setting. An acquisition never alters the
exchange process. The scanner must be available in the corresponding
window and cannot replay a missed one. A conditioning variable such
as a known job queue can affect offered scan settings, provided its
law is fixed across the two modes and it supplies no extra channel
about exchange identities.

Historical runs already owe inspection service that occupies a
scanner for at least one of those windows. Approval of new training
runs may instead reserve the diagnostic scanner throughout the
relevant windows by declining an optional job \emph{before} accepting
it. Count its foregone value \(\kappa\); do not cancel live inspection
work by assuming it has been settled. Historical status and safe
fallback outcomes do not encode the mode; in particular there are
no persistent observed identities at both ends, exact end-to-end
tracking, or past randomized full preparatory trials with observed
success. If any such measurement is feasible it is a different,
possibly better, matched experiment. The claims below concern one
new diagnostic run for a future preparation, preceded by any number
of these historical runs.

\section{Two indistinguishable historical laws}
Any strict subset of the exchange coordinates is jointly uniform
under either source law. For an adaptive transcript ending without
reading all stages, condition on the controller's seed and fixed
channel options along that transcript. The event of taking the
specified actions is implied by obtaining its specified readouts.
Its likelihood is therefore the product of fair readout probabilities
on the scanned subset, under either law. This covers the decision
to skip a window based on earlier readings, provided feasibility is
decided from the declared local information and initial capacities.
No adaptive selection of future historical runs changes this: the
likelihood of each new incomplete transcript given the prior data
is the same, so induct over runs. Even exact records over different
proper subsets in different runs cannot recover a same-run parity.

The sharp worst-law error with those data alone is \(1/2\).
The average error under equal priors is \(1/2\) for any decision
because the distributions coincide. Randomizing the two placements
equally attains this bound in each law. This is not failure to
estimate the separate \(X_i\) marginals: those marginals, and even
every proper joint marginal of \(X\), are exactly known in the model.

\section{One diagnostic run: the exact frontier}

For a complete fixed sequence of settings \(u\), set
\(r(u)=\prod_i(1-2p_i(u_i))\). Set \(r=0\) for an incomplete schedule.
With all readings the two likelihoods at \(y\) are
\[
 P_s(y\mid u)=2^{-n}\bigl[1+(-1)^{s+\oplus_i y_i}r(u)\bigr].
\]
To see this, condition on the noise vector. The specified readout
then has probability \(2^{-(n-1)}\) or zero depending on its parity
after subtracting that vector. Independence gives
\(\mathbb E[(-1)^{\oplus_i Z_i}]=\prod_i(1-2p_i(u_i))\).
For adaptive setting selection the same leaf formula holds
with the settings on that leaf: the rule selects them whenever
those exact readouts occur. Every strict-prefix likelihood is
still fair. In particular the two continuations after the last
setting is chosen occur with combined probability equal to their
unbiased mass under either source law. An incomplete leaf is
unbiased on its entire transcript.

This yields two conclusions. First, under an adaptive policy let
\(\nu\) be the distribution of acquisition schedules in a reference
process in which every emitted bit is independent and fair.
Actual schedule probabilities are also \(\nu\) in each law:
if complete, sum over the final bit, since the setting for that
bit is chosen before it is seen; otherwise all acquired coordinates
form a proper subset. Consequently both the expected incurred
cost \(\mathbb E_\nu C(u)\) and the exact variation distance
\(\mathbb E_\nu r(u)\) are schedule averages. Here \(C\) includes
the access cost if paid, plus the actual acquisition charges.
A complete readout recommends its own parity. Conditional on
the strict prefix the error of this recommendation is
\((1-r(u))/2\) in either law. An incomplete transcript induces
a fair guess. Thus the minimum worst-law error with this
experiment policy is exactly
\[
         \frac{1-\mathbb E_\nu r(u)}2.
\]
For the lower bound, the equal-prior best testing error on a common
transcript space is one half of the sum of the pointwise minima
of the two probabilities, or half of one minus their variation
distance. The proposed rule attains that bound in both laws.

Second, assume capacity interactions depend only on the actual
windows, selected actions and their charges, not on the values
reported by scans; fix job-queue conditions when determining that
envelope. Each schedule in \(\nu\)'s support is then feasible if
chosen in advance. Draw the schedule with distribution \(\nu\)
before the run and apply the fixed-schedule decisions above:
it has the same cost and error. Conversely, randomly reserving
any feasible schedules gives their corresponding averages. The
best attainable error under an expected cost cap is therefore
the lower boundary of the convex hull of the fixed-schedule pairs
\((C(u),(1-r(u))/2)\). For minimum expected total cost with unit
price per resource and loss \(\ell\) for a wrong forecast,
randomization cannot beat the best fixed schedule:
minimize \(C(u)+\ell(1-r(u))/2\). Pay to clear the old scan
window only if some newly available complete schedule has
\(C(u)<\ell r(u)/2\); the uninformative zero-cost alternative has
cost \(\ell/2\).

For example, with two windows, scans that each cost \(1\) and
err independently with probability \(0.1\) give testing error
\(0.18\); cheap scans at \(1/3\) each err with probability \(0.25\).
With \(\ell=10\), the best complete schedule uses two accurate
scans and costs \(3.8+\kappa\) including service loss; the mixed
and cheap schedules cost \(13/3+\kappa\) and \(53/12+\kappa\).
A prepared run is worth arranging precisely when \(\kappa<1.2\)
(indifference at equality), assuming no cheaper outside experiment.
The access cost is sunk even when the diagnosis is wrong, and
a wrong placement still incurs the slower fallback.

\section{Failure of the reduction and comparison}

The independence qualification is material. For two \(10\%\)-error
readings, independent errors make the net parity wrong with
probability \(18\%\). If the same calibration slip flips both,
the marginal errors are unchanged and the parity is always right.
With only the two 10\% marginals, the sharp net error is anywhere
between 0 and 20\%: it equals 0.2 less twice the probability of
a double error, which can range from 0 to 0.1. At access cost
1.1 the two accurate scans are worth arranging under independence
(total 4.9 versus 5 without scans), but not for those same two scans
under the worst case allowed by their separate calibrations (5.1).
The same separate calibration reports therefore do not justify
the same release decision. The needed calibration concerns the
\emph{paired} experiment; time spent reserving read windows does
not by itself supply that calibration. A report that changes
future acquisition feasibility or triggers an informative physical
action also falls outside the open-loop reduction. A direct
identity observation available at the right time bypasses the
parity experiment entirely.

The ordinary likelihood-ratio test on the full admissible
transcript attains the same result. This is a boundary for what
changes when live commitments constrain the experiment, not an
improved classifier. \href{https://arxiv.org/abs/1205.0858}{Nitinawarat,
Atia and Veeravalli} treat controlled hypothesis testing and
derive asymptotic comparisons of sensing policies; this benchmark
calculates exact finite-run costs for a dependent latent source and
explicit window commitments. The contract's
\href{https://proceedings.mlr.press/v286/marmarelis25a.html}{off-policy
causal sensitivity baseline} studies hidden confounding and policy
evaluation. The present counterexample has two stated laws and a
known observation mechanism; it neither reproduces that estimator
nor establishes a broadly valid method for identifying forecasts
from observational field data. Exact enumerations in the companion
code compare adaptive transcript calculations with complete source
histories. This is partial R3b evidence, not a novelty or acceptance
claim for the whole obligation.
\end{document}

